\section{The problem and the origin of the construction}
\label{sec:introduction}

The objective is a nontrivial four-dimensional quantum Yang--Mills theory on
\(\RR^4\), initially for gauge group \(SU(2)\). The intended construction
starts with the positive Wilson lattice measure, removes the lattice cutoff
along a controlled weak-coupling trajectory, and constructs local
gauge-invariant quantum fields. The problem is to carry the \emph{actual}
measure, its normalization, and its observables through these limits. A
Gaussian comparison, an effective action on a small-field chart, or a
finite-dimensional space of ultraviolet symbols is not by itself that
construction.

This work is a continuation of the public project of Levent Alpoge and
Ethan Sussman, \emph{Yang--Mills LLM experiment} \cite{Upstream}. Their
deposited material develops Ba{\l}aban's renormalization-group programme,
action-density source estimates, a scalar reconstruction route, and a
reflection-positivity approach to continuous rotations. Those starting
ideas belong to that project and to the cited constructive field theory
literature. The examined snapshot is the commit
\begin{center}\small\path{e392ef45013a581664bab76b90867170b30fe10e}.\end{center}
Its README reports an apparent argument on \(\RR^4\); we do not recast it
as merely a finite-volume proposal. The displayed hypotheses and the
narrower observable conclusions nevertheless have to be checked
individually. In particular, a public existence claim does not discharge
an unverified domain or observable estimate in the continuation.

The continuation develops explicit local constructions, measure-matching
identities, source transport estimates, and algebraic tests of whether a
scalar construction can recover the Yang--Mills observable algebra. The
latest improvements concern local adjusted-average right inverses,
physical currents, complex-domain margins, dated source transport, and
the oriented cubic curvature sector. The paper presents these as
mathematical statements with their premises. It does not assert a
completed continuum Yang--Mills existence theorem.

\subsection{The precise target}

Write \(\mathfrak g=\mathfrak{su}(2)\). A successful endpoint includes a
Hilbert space \(\HH\), a vacuum \(\Omega\), a strongly continuous
positive-energy representation of the Poincar\'e group, and a compatible
family of local operator-valued tempered distributions corresponding to
gauge-invariant polynomials in the curvature and its covariant
derivatives. The fields must have the prescribed tensor transformation
laws, adjoints, derivative identities, and Yang--Mills short-distance
normalization. A nonzero positive-energy vacuum excitation supplies a
nontriviality witness. This fixes the meaning of existence used here;
it is stronger than constructing a scalar field whose two-point function
resembles a free curvature square. The relation to the customary
Yang--Mills problem is discussed in \cite{JaffeWitten}.

We distinguish three tasks. First, construct a nontrivial scalar sector
from the action density. Second, construct and identify the full local
observable hierarchy in that sector or in a verified extension of it.
Third, prove a spectral gap. The third task is separate and is not solved
by the results below. Nor do the special \(SU(2)\) invariant-theory
arguments establish the construction for every compact simple group.

\subsection{Lattice conventions and normalization}

On a finite oriented periodic lattice \(\Lambda\) of spacing \(\eta\),
assign \(U_e\in SU(2)\) to each positively oriented edge, with
\(U_{\bar e}=U_e^{-1}\). Haar measure is normalized to one. For an
unoriented plaquette set \(\mathcal P\), counted once,
\begin{equation}
 A_\Lambda(U)=\sum_{p\in\mathcal P}
       (1-\operatorname{Re}\tr_N U(\partial p)),\qquad
 d\mu_{\Lambda,g}=Z_{\Lambda,g}^{-1}e^{-g^{-2}A_\Lambda(U)}\prod_e dU_e,
 \quad \tr_N 1=1. \label{eq:wilson}
\end{equation}
Numerical beta coefficients depend on this choice, on plaquette counting,
and on the normalization of the continuum Maxwell form. Every comparison
below retains its stated convention. In the cubic algebra it is convenient
instead to use Hermitian generators \(t_a=\sigma_a/2\) and the ordinary
matrix trace \(\Tr(t_at_b)=\delta_{ab}/2\); then
\(\tr_N=\Tr/2\). These traces must not be silently interchanged. With
this Wilson plaquette convention the standard bare canonical coupling
satisfies \(g_{\rm can}=2g_{\rm W}\), equivalently
\(\beta=g_{\rm W}^{-2}=4g_{\rm can}^{-2}\). The later quoted
\(8g/Q^2\) cubic coefficient uses \(g=g_{\rm can}\), not
\(g_{\rm W}\). Its application also requires the running source's field
normalization. A local section's declared coupling convention therefore
governs its numerical coefficients.

For local lattice action-density observables \(O_i\), define
\begin{equation}
 Z(z)=\int e^{-g^{-2}A_\Lambda(U)+\sum_i z_iO_i(U)}dU,
 \quad \partial_{z_1}\partial_{z_2}\log Z(0)
       =\Cov_{\mu_{\Lambda,g}}(O_1,O_2). \label{eq:sources}
\end{equation}
This identity fixes the normalization of marked descendants. Source
derivatives act on the full expression, including extracted scalar terms
and the partition denominator. A vacuum contribution which is independent
of the field may still depend on \(z\); its derivatives cannot be deleted.

We use \(g\) for the running coupling, \(\gamma=g^{-2}\) for its inverse
square when that coordinate is appropriate, \(b\) for an analytic
background radius, and \(T\) for a retained Gaussian fluctuation window.
The commonly used inner radius is \(b=cg\log\gamma\), with
\(g\ll b\), and an inner window may be chosen
\(T=(\log\gamma)^t\), \(1/2<t<1\). Outer preparation radii and windows
are different quantities. A fixed blocking factor is denoted by \(L\);
\(q\) counts a fixed number of primitive steps. Constants depending on
\(L,q\) are allowed; dependence on ultraviolet depth, volume, or birth age
is not allowed in a claim of uniformity. A scale ratio \(h\leq1\) is
always identified at its use. It is not an additional small coupling.

\subsection{Status conventions}

There are three kinds of statements in this manuscript.
\begin{enumerate}
 \item An \emph{exact statement} is a finite-dimensional identity,
 an elementary analytic implication, or an algebraic calculation with
 all its mathematical hypotheses stated here.
 \item A \emph{source-relative statement} uses a specified estimate for
 Ba{\l}aban's maps, or a regular-kernel moment/locality package. Its proof
 establishes the stated implication, but does not reprove the imported
 package on every mixed large-field history.
 \item A \emph{conditional continuum statement} specifies the missing
 uniform estimates and proves what follows from them. Its conclusion
 is not an established property of the Wilson continuum limit until
 those premises are verified for the actual recursion.
\end{enumerate}
These distinctions concern logical scope, not confidence labels. A proof
of an implication is useful only when its antecedent is kept visible.
Repeated checks of a notebook do not replace an independent verification
of a source hypothesis. The companion records counterexamples and
superseded statements; the present paper uses the final audited scope.

\section{Architecture of the existence argument}
\label{sec:architecture}

\subsection{The four connected constructions}

The first construction is local and variational. One controls constrained
minima, their gauge representatives, physical currents, Gaussian
fluctuations, and the complete reduced measure. It supplies a well-defined
homogeneous renormalization step and identifies its leading running
coefficient. The local calculation is useful only after its constraint,
gauge, determinant, cutoff, and observable conventions match the
corresponding source integral.

The second construction is a full-history source estimate. A localized
insertion is born at a specified scale and is transported through every
later regular and large-field operation. Four-dimensional root counting
costs \(h^{-4}\); a curvature-square insertion has neutral scaling
\(h^4\). Their product is only order one. Improvement therefore comes
from a verified cancellation, a positive extra power of \(h\), or the
actual smallness of its birth. It cannot be obtained by forgetting old
roots or replacing a birth date by a later scale.

The third construction converts source estimates into scalar continuum
objects. Positivity is taken from the original Wilson measure. Copying
and reflection arguments must use one normalized denominator and one
state. Continuum rotations, uniform growth in the number of insertions,
and the passage from Euclidean forms to operators are separate steps.

The fourth construction enlarges the observable family. Scalar fusion
has enough Gaussian quadratic symbols to span all 21 curvature-pair
directions. It cannot, at Gaussian order, produce the color-oriented
dimension-six sector. An interaction-sensitive cubic calculation and
an actual operator remainder estimate are required. Even the 21-field
construction needs continuous tensor covariance and a relative Gram
bound in the limiting theory.

\subsection{What is currently proved, and what is still a premise}

\begin{longtable}{@{}>{\raggedright\arraybackslash}p{.23\linewidth}>{\raggedright\arraybackslash}p{.40\linewidth}>{\raggedright\arraybackslash}p{.29\linewidth}@{}}
\toprule
Part & Supported conclusion & Remaining attachment \\
\midrule\endhead
Local regular recursion & Constraint/measure identities, local analytic
response, homogeneous coupling and drift reductions & Uniform matching
to the complete mixed-history density \\
Physical current & Local adjusted-average inverse and current bound in
one admitted determining hierarchy & Admission of every modified
intrinsic base on its complete collar \\
Dated source calculus & Conditional summation, phase cancellation,
copying and normalized-pressure implications & Full all-spectator
first-contact estimate for the literal histories \\
Scalar reconstruction & Positivity, growth, slab and operator reductions
under their stated source/rotation inputs & Application to one controlled
continuum hierarchy and a nontrivial vacuum phase \\
Quadratic and cubic symbols & Exact 21-dimensional span; oriented cubic
algebra; fixed-\(q\) parity and first-order expansion & Actual relative
Gram, long-prefix cubic transport and uniform remainders \\
Full observable theory & A sufficient common-Hilbert-space completion
criterion & Tensor covariance, higher jets, OPE and stress-tensor
normalization \\
\bottomrule
\end{longtable}

The sharpest current source question is formulated in
Section~\ref{sec:admission}. It concerns preparation-changed datum slots,
not an unspecified need for ``better estimates.'' Solving it would unlock
an important branch of the source argument, but it would not automatically
prove the full observable hierarchy or continuous tensor covariance.

\subsection{Conventions for the accompanying record}

The main paper consolidates the results that participate in this route.
Proofs of elementary reductions are included; longer local constructions
are also preserved in the numbered notebooks. The accompanying archive
contains the mathematical research notes, proof fragments, correction
records, and computational checks, with relative paths and SHA-256
hashes. Filenames quoted below are locators into that archive. The
research companion contains the parked branches and negative tests.
The original upstream papers are identified by version and links rather
than silently merged into the continuation's authorship or claims.
